\DocumentMetadata{
  pdfversion=2.0,
  pdfstandard=ua-2,
  lang=en-US,
  tagging=on,
  tagging-setup={math/setup=mathml-SE,tabsorder=structure}
}
\documentclass[11pt,letterpaper]{article}
\usepackage[margin=0.68in,includefoot]{geometry}
\usepackage{newcomputermodern}
\usepackage{amsmath,amscd,mathtools}
\usepackage{tikz,adjustbox}
\providecommand{\square}{\mdlgwhtsquare}
\usepackage[hidelinks]{hyperref}
\hypersetup{
  pdftitle={Analysis PhD Exam, August 2024},
  pdfauthor={University of Florida Department of Mathematics},
  pdfsubject={Graduate mathematics examination},
  pdfdisplaydoctitle=true,
  bookmarksopen=true
}
\setcounter{secnumdepth}{0}
\setlength{\parindent}{0pt}
\setlength{\parskip}{0.28em}
\setlength{\emergencystretch}{3em}
\setlength{\leftmargini}{2.15em}
\setlength{\leftmarginii}{2.65em}
\setlength{\leftmarginiii}{3em}
\pagestyle{empty}
\raggedbottom
\allowdisplaybreaks
\NewTaggingSocketPlug{block/list/label}{examproblem}{%
  \tagstructbegin{tag=Lbl}\tagstructend%
  \tagstructbegin{tag=\LBody}%
  \tagstructbegin{tag=\ProblemHeadingTag,title-o={\ProblemHeadingText}}%
  \tagmcbegin{tag=\ProblemHeadingTag,actualtext={\ProblemHeadingText}}%
  #2%
  \tagmcend%
  \tagstructend%
}
\newenvironment{examproblems}[2]{%
  \def\ProblemHeadingTag{#2}%
  \AssignTaggingSocketPlug{block/list/label}{examproblem}%
  \begin{enumerate}%
  \setcounter{enumi}{#1}%
  \renewcommand{\labelenumi}{\textbf{\theenumi.}}%
  \setlength{\topsep}{0.35em}%
  \setlength{\partopsep}{0pt}%
  \setlength{\itemsep}{0.48em}%
  \setlength{\parsep}{0.18em}%
}{\end{enumerate}}
\newenvironment{examsubparts}{%
  \AssignTaggingSocketPlug{block/list/label}{default}%
  \begin{enumerate}%
  \setlength{\topsep}{0.18em}%
  \setlength{\partopsep}{0pt}%
  \setlength{\itemsep}{0.16em}%
  \setlength{\parsep}{0.1em}%
}{\end{enumerate}}
\newenvironment{examinstructions}{%
  \par\begingroup\itshape
}{%
  \par\endgroup\vspace{0.18em}
}
\makeatletter
\renewcommand\subsection{\@startsection{subsection}{2}{0pt}%
  {-1.25ex plus -.25ex minus -.1ex}%
  {0.55ex plus .1ex}%
  {\normalfont\large\bfseries}}
\renewcommand{\@maketitle}{%
  \newpage\null\vskip -1em%
  {\footnotesize\bfseries
    \href{https://www.ufl.edu/}{University of Florida}, \href{https://math.ufl.edu/}{Department of Mathematics}\par}%
  \vskip -0.5em%
  \tagstructbegin{tag=H1,title={Analysis PhD Exam, August 2024}}%
  \tagpdfparaOff%
  \tagmcbegin{tag=H1}%
    {\LARGE\bfseries\href{https://gma.math.ufl.edu/exams/analysis-phd/}{Analysis PhD Exam}, August 2024\par}%
  \tagmcend%
  \tagpdfparaOn%
  \tagstructend%
  \vskip 0.25em%
  \tagmcbegin{artifact}%
    \hrule height 1pt%
  \tagmcend%
  \vskip 1em%
}
\makeatother
\title{Analysis PhD Exam, August 2024}
\author{}
\date{}
\begin{document}
\maketitle
\thispagestyle{empty}
\begin{examinstructions}
Write solutions in a neat and logical fashion, giving complete reasons for all steps. Attempt SIX problems.
\end{examinstructions}
\begin{examproblems}{0}{H2}
\def\ProblemHeadingText{Problem 1.}
\item
For any two of the following, either give an example or explain why no example exists.
\begin{examsubparts}
\item[\textnormal{a)}]
a sequence of Lebesgue measurable functions on \(\mathbb R\) which converges in measure but not Lebesgue a.e.
\item[\textnormal{b)}]
a closed set \(E\subset[0,1]\) with positive Lebesgue measure, but which contains no open intervals
\item[\textnormal{c)}]
a decreasing sequence of nonnegative measurable functions on \(\mathbb R\) such that \(\lim \int f_n \ne \int \lim f_n\)
\end{examsubparts}
\def\ProblemHeadingText{Problem 2.}
\item
State Tonelli’s theorem, and give an example to show that the~\(\sigma\)-finiteness hypothesis is necessary.
\def\ProblemHeadingText{Problem 3.}
\item
This problem has two parts.
\begin{examsubparts}
\item[\textnormal{a)}]
State the Lebesgue–Radon–Nikodym theorem.
\item[\textnormal{b)}]
State and prove a version of the chain rule for Radon–Nikodym derivatives.
\end{examsubparts}
\def\ProblemHeadingText{Problem 4.}
\item
Suppose that \(f:\mathbb R\to\mathbb C\) is a Lebesgue measurable function with the property that \(fg\in L^1(\mathbb R)\) for every \(g\in L^2(\mathbb R)\). Must~\(f\) belong to \(L^2(\mathbb R)\)? Prove, or give a counterexample.
\def\ProblemHeadingText{Problem 5.}
\item
This problem has two parts.
\begin{examsubparts}
\item[\textnormal{a)}]
Show that for every integer \(n\ge 1\), there exists a function \(f_n\in L^2[0,1]\) such that for every polynomial~\(p\) of degree at most~\(n\),
\[
p(0)=\int_0^1 p(x)f_n(x)\,dx.
\]
\item[\textnormal{b)}]
Prove that there is no \(L^2\) function~\(f\) such that the above holds for all polynomials~\(p\) (independent of the degree~\(n\)).
\end{examsubparts}
\def\ProblemHeadingText{Problem 6.}
\item
State the Baire category theorem. Prove that if \(\mathcal X\) is an infinite-dimensional Banach space, then \(\mathcal X\) does not have a countable basis.
\def\ProblemHeadingText{Problem 7.}
\item
Let \(\mathcal X\) be a Banach space and \(\mathcal Y\) a normed vector space. Let \((T_n)\) be a sequence of bounded linear operators from \(\mathcal X\) to \(\mathcal Y\). Suppose that \(\lim_n T_nx\) exists for each \(x\in\mathcal X\). Prove that the linear transformation \(Tx:=\lim_n T_nx\) is bounded from \(\mathcal X\) to \(\mathcal Y\).
\def\ProblemHeadingText{Problem 8.}
\item
For each of the following statements, determine if it is true or false, and sketch a proof of your claim.
\begin{examsubparts}
\item[\textnormal{a)}]
If~\(\nu\) is a signed measure and \(E,F\) are measurable sets with \(\nu(E)\ge 0\) and \(\nu(F)\ge 0\), then \(\nu(E\cup F)\ge 0\).
\item[\textnormal{b)}]
If \((X,\mathcal M)\) and \((Y,\mathcal N)\) are measurable spaces and \(E\subset X\), \(F\subset Y\) are sets such that \(E\times F\) is \(\mathcal M\otimes\mathcal N\)-measurable, then \(E\in\mathcal M\) and \(F\in\mathcal N\).
\item[\textnormal{c)}]
\(L^1(\mathbb R)\) is reflexive.
\end{examsubparts}
\end{examproblems}
\end{document}
